% ---------------------------------------------------------------------
\chapter{Die Befehle im Einzelnen}
\label{kap:befehle}

\begin{lernziele}
\begin{itemize}[nosep]
  \item für jeden Alltagsbefehl wissen, welchen Weg im mentalen Modell er geht,
  \item jede Option eines Befehls einzeln verstehen und gezielt kombinieren,
  \item die Ausgaben von \cmd{status}, \cmd{diff}, \cmd{log} und \cmd{push} lesen können.
\end{itemize}
\end{lernziele}

Die Befehle sind in der Reihenfolge des Arbeitszyklus geordnet. Jeder Abschnitt beginnt mit
dem Zweck, zerlegt dann den Befehl in seine Teile und schliesst mit Hinweisen aus der Praxis.

% =====================================================================
\befehl{git status}{den Zustand abfragen}

\textbf{Zweck:} zeigt, welche Änderungen sich in welchem der vier Orte befinden. Der Befehl
verändert nichts. Er ist der meistverwendete Git-Befehl überhaupt.

\begin{lstlisting}[style=befehle]
git status
\end{lstlisting}

\begin{lstlisting}[style=ausgabe]
On branch main
Your branch is up to date with 'origin/main'.

Changes to be committed:
        modified:   python/uebung01/al/aufgabe03/recursive_sum.py

Changes not staged for commit:
        modified:   README.md

Untracked files:
        notizen.txt
\end{lstlisting}

Die Ausgabe gliedert sich nach den Orten des mentalen Modells:

\begin{center}
\renewcommand{\arraystretch}{1.3}
\begin{tabularx}{\linewidth}{@{}>{\raggedright\arraybackslash}p{4.6cm}>{\raggedright\arraybackslash}X@{}}
  \toprule
  \textbf{Abschnitt} & \textbf{Bedeutung} \\
  \midrule
  Changes to be committed      & in der Staging Area, kommt in den nächsten Commit (grün) \\
  Changes not staged for commit& im Arbeitsordner geändert, aber nicht vorgemerkt (rot) \\
  Untracked files              & neue Dateien, die Git noch nie gesehen hat (rot) \\
  up to date with origin/main  & gleicher Stand wie beim letzten Abgleich mit GitHub \\
  \bottomrule
\end{tabularx}
\end{center}

\subsection*{Kurzform}

\begin{lstlisting}[style=befehle]
git status -sb
\end{lstlisting}

\begin{teile}
  -s  & \textbf{s}hort: eine Zeile pro Datei statt ganzer Sätze \\
  -b  & \textbf{b}ranch: zeigt trotzdem die Zeile mit dem Branch \\
\end{teile}

\begin{lstlisting}[style=ausgabe]
## main...origin/main [ahead 1]
M  python/uebung01/al/aufgabe03/recursive_sum.py
 M README.md
?? notizen.txt
\end{lstlisting}

Die zwei Zeichen vor jedem Dateinamen sind zwei Spalten: Die \textbf{linke} beschreibt die
Staging Area, die \textbf{rechte} den Arbeitsordner. \cmd{[ahead 1]} bedeutet: ein Commit
mehr als auf GitHub.

\begin{center}
\renewcommand{\arraystretch}{1.25}
\begin{tabularx}{\linewidth}{@{}>{\ttfamily}l>{\raggedright\arraybackslash}X@{}}
  \toprule
  \normalfont\textbf{Zeichen} & \textbf{Bedeutung} \\
  \midrule
  M\textvisiblespace  & geändert und vorgemerkt \\
  \textvisiblespace M & geändert, nicht vorgemerkt \\
  MM                  & vorgemerkt und danach erneut geändert \\
  A\textvisiblespace  & neue Datei, vorgemerkt (\emph{added}) \\
  D\textvisiblespace  & gelöscht, vorgemerkt (\emph{deleted}) \\
  ??                  & unbekannte Datei (\emph{untracked}) \\
  \bottomrule
\end{tabularx}
\end{center}

\begin{erfahrung}
Vor jedem Commit \cmd{git status}. Die zwei Sekunden verhindern fast jeden Fall von „das
hätte nicht mit sollen“.
\end{erfahrung}

% =====================================================================
\befehl{git add}{Änderungen vormerken}

\textbf{Zweck:} legt Änderungen aus dem Arbeitsordner in die Staging Area.
\textbf{Weg:} Arbeitsordner $\rightarrow$ Staging Area.

\begin{lstlisting}[style=befehle]
git add python java
\end{lstlisting}

Die Argumente sind Dateien oder Ordner, \textbf{relativ zum aktuellen Ordner}. Bei einem
Ordner werden alle Änderungen darin und in allen Unterordnern vorgemerkt.

\begin{center}
\renewcommand{\arraystretch}{1.25}
\begin{tabularx}{\linewidth}{@{}>{\ttfamily}l>{\raggedright\arraybackslash}X@{}}
  \toprule
  \normalfont\textbf{Aufruf} & \textbf{Was vorgemerkt wird} \\
  \midrule
  git add datei.py & genau diese Datei \\
  git add python   & alles unter \cmd{python/} \\
  git add .        & alles im aktuellen Ordner und darunter; der Punkt bedeutet „hier“ \\
  git add -A       & alle Änderungen im ganzen Repository (\emph{all}), auch gelöschte Dateien \\
  git add -p       & interaktiv Stück für Stück (\emph{patch}); Git fragt bei jedem Abschnitt \cmd{y} oder \cmd{n} \\
  \bottomrule
\end{tabularx}
\end{center}

Vorgemerkt werden nur \textbf{Änderungen}. Unveränderte Dateien sind bereits im letzten
Commit enthalten und werden nicht erneut gespeichert.

\begin{praxis}
Erfahrene Entwicklerinnen und Entwickler verwenden \cmd{git add -A} zurückhaltend. Sobald
mehrere Anliegen gleichzeitig offen sind, merken sie gezielt Pfade vor oder nutzen
\cmd{git add -p}, um sogar innerhalb einer Datei zu trennen. So bleibt jeder Commit auf ein
Anliegen beschränkt, was das Prüfen und spätere Rückgängigmachen erleichtert.
\end{praxis}

% =====================================================================
\befehl{git diff}{Unterschiede anzeigen}

\textbf{Zweck:} zeigt Zeile für Zeile, was sich zwischen zwei Ständen unterscheidet
(\emph{diff} von \emph{difference}). Welche zwei Stände verglichen werden, hängt von den
Optionen ab:

\begin{center}
\renewcommand{\arraystretch}{1.3}
\begin{tabularx}{\linewidth}{@{}>{\ttfamily}l>{\raggedright\arraybackslash}p{4.4cm}>{\raggedright\arraybackslash}X@{}}
  \toprule
  \normalfont\textbf{Aufruf} & \textbf{vergleicht} & \textbf{beantwortet die Frage} \\
  \midrule
  git diff          & Arbeitsordner $\leftrightarrow$ Staging Area & Was habe ich geändert, aber noch nicht vorgemerkt? \\
  git diff --cached & Staging Area $\leftrightarrow$ letzter Commit & Was kommt in den nächsten Commit? \\
  git diff HEAD     & Arbeitsordner $\leftrightarrow$ letzter Commit & Was ist insgesamt seit dem letzten Commit anders? \\
  \bottomrule
\end{tabularx}
\end{center}

\begin{hintergrund}
\cmd{--cached} heisst so, weil die Staging Area in Git intern früher \emph{cache} genannt
wurde. Da der Name verwirrend ist, gibt es seit einigen Jahren das gleichbedeutende
\cmd{--staged}.
\end{hintergrund}

\subsection*{Die Ausgabe lesen}

\begin{lstlisting}[style=ausgabe]
--- a/recursive_sum.py
+++ b/recursive_sum.py
@@ -3,4 +3,6 @@
 def recursive_sum(n):
-    pass
+    if n == 0:
+        return 0
+    return n + recursive_sum(n - 1)
\end{lstlisting}

\begin{itemize}
  \item \cmd{---} kennzeichnet die alte Fassung (a), \cmd{+++} die neue (b).
  \item \cmd{@@ -3,4 +3,6 @@} bedeutet: Der Ausschnitt beginnt in Zeile 3; vorher umfasste er
        4 Zeilen, jetzt 6.
  \item Zeilen mit \cmd{-} wurden entfernt, Zeilen mit \cmd{+} hinzugefügt. Zeilen ohne
        Zeichen sind unverändert und dienen der Orientierung.
  \item Lange Ausgaben zeigt Git seitenweise: Leertaste blättert weiter, \cmd{q} beendet.
\end{itemize}

\subsection*{Nur die Übersicht: \texorpdfstring{\cmd{--stat}}{--stat}}

\begin{lstlisting}[style=befehle]
git diff --cached --stat
\end{lstlisting}

\begin{lstlisting}[style=ausgabe]
 java/uebung01/as/aufgabe03/RecursiveSum.java  | 12 +++++++---
 python/uebung01/al/aufgabe03/recursive_sum.py |  5 +++--
 2 files changed, 13 insertions(+), 4 deletions(-)
\end{lstlisting}

\cmd{--stat} (\emph{statistics}) zeigt pro Datei nur die Anzahl geänderter Zeilen und das
Verhältnis von hinzugefügten (\cmd{+}) zu entfernten (\cmd{-}) Zeilen.

Die beiden Optionen sind voneinander unabhängig: \cmd{--cached} bestimmt, \textbf{was}
verglichen wird, \cmd{--stat}, \textbf{wie} es dargestellt wird. Am besten probiert man sie
einzeln aus:

\begin{center}
\renewcommand{\arraystretch}{1.25}
\begin{tabularx}{\linewidth}{@{}>{\ttfamily}l>{\raggedright\arraybackslash}X@{}}
  \toprule
  \normalfont\textbf{Aufruf} & \textbf{Ergebnis} \\
  \midrule
  git diff                      & Details, nur nicht Vorgemerktes \\
  git diff --stat               & Übersicht, nur nicht Vorgemerktes \\
  git diff --cached             & Details, nur Vorgemerktes \\
  git diff --cached --stat      & Übersicht, nur Vorgemerktes \\
  git diff --cached --name-only & nur die Dateinamen \\
  \bottomrule
\end{tabularx}
\end{center}

% =====================================================================
\befehl{git restore}{Änderungen zurücknehmen}

\textbf{Zweck:} holt Dateien aus der Staging Area zurück oder verwirft Änderungen.
\textbf{Weg:} Staging Area $\rightarrow$ Arbeitsordner.

\begin{lstlisting}[style=befehle]
git restore --staged datei.py          # aus der Vormerkung nehmen, Datei bleibt
git restore datei.py                   # Änderungen an der Datei verwerfen
\end{lstlisting}

\begin{teile}
  --staged & nur die Vormerkung aufheben; der Inhalt der Datei bleibt unverändert \\
  (ohne)   & die Datei im Arbeitsordner auf den letzten Commit zurücksetzen \\
\end{teile}

\begin{achtung}
\cmd{git restore datei.py} \emph{ohne} \cmd{--staged} löscht alle nicht committeten
Änderungen an dieser Datei endgültig. Es gibt keinen Papierkorb. Vorher mit
\cmd{git diff datei.py} prüfen, was verloren ginge.
\end{achtung}

% =====================================================================
\befehl{git commit}{einen Zwischenstand speichern}

\textbf{Zweck:} speichert alles Vorgemerkte als neuen Commit.
\textbf{Weg:} Staging Area $\rightarrow$ lokales Repository.

\begin{lstlisting}[style=befehle]
git commit -m "ADS U01: rekursive Summe implementiert"
\end{lstlisting}

\begin{teile}
  -m & \textbf{m}essage: die Beschreibung direkt angeben; ohne \cmd{-m} öffnet sich ein Texteditor \\
\end{teile}

Der Commit existiert danach nur \textbf{lokal}. GitHub erfährt erst beim Push davon. Wie man
gute Nachrichten schreibt, behandelt Kapitel~\ref{kap:commits}.

% =====================================================================
\befehl{git log}{die Geschichte ansehen}

\textbf{Zweck:} listet die Commits auf, den neuesten zuerst.

\begin{lstlisting}[style=befehle]
git log --oneline --graph --all -10
\end{lstlisting}

\begin{teile}
  --oneline & eine Zeile pro Commit: kurzer Hash und erste Zeile der Nachricht \\
  --graph   & zeichnet die Verzweigungen der Branches als Linien \\
  --all     & alle Branches statt nur des aktuellen \\
  -10       & nur die letzten zehn Commits \\
\end{teile}

\begin{lstlisting}[style=ausgabe]
*   3b1a03a Merge pull request #1 from gladtobethere/ads
|\
| * db753be ADS: Uebungen U01-U06
|/
* 19ce440 ADS-Uebungsordner anlegen
\end{lstlisting}

Die Grafik zeigt: Auf einem Branch entstand \cmd{db753be}, der danach mit einem
Merge-Commit (\cmd{3b1a03a}) in \cmd{main} übernommen wurde.

% =====================================================================
\befehl{git switch}{Branch wechseln oder anlegen}

\textbf{Zweck:} wechselt den aktuellen Branch, HEAD wandert mit.

\begin{lstlisting}[style=befehle]
git switch main                        # auf bestehenden Branch wechseln
git switch -c ads                      # neuen Branch anlegen und wechseln
\end{lstlisting}

\begin{teile}
  -c & \textbf{c}reate: den Branch zuerst anlegen, ausgehend vom aktuellen Commit \\
\end{teile}

Nicht committete Änderungen wandern beim Anlegen eines neuen Branches mit. Wer versehentlich
auf \cmd{main} gearbeitet hat, rettet die Arbeit deshalb einfach mit \cmd{git switch -c <name>}.

\begin{hintergrund}
\cmd{git switch} gibt es seit 2019. Ältere Anleitungen verwenden \cmd{git checkout -b ads}.
Das bewirkt dasselbe; \cmd{checkout} wurde aufgeteilt, weil es zu viele verschiedene
Aufgaben hatte.
\end{hintergrund}

% =====================================================================
\befehl{git branch}{Branches auflisten und löschen}

\begin{lstlisting}[style=befehle]
git branch                             # lokale Branches; * markiert den aktuellen
git branch -a                          # auch die Branches auf GitHub (all)
git branch -d ads                      # löschen (delete), nur wenn bereits übernommen
\end{lstlisting}

\cmd{-d} ist bewusst vorsichtig und verweigert das Löschen eines Branches, dessen Arbeit
noch nicht übernommen wurde. Das grosse \cmd{-D} erzwingt es und kann Arbeit kosten.

% =====================================================================
\befehl{git push}{Commits hochladen}

\textbf{Zweck:} schickt lokale Commits zum Remote.
\textbf{Weg:} lokales Repository $\rightarrow$ GitHub.

\begin{lstlisting}[style=befehle]
git push -u origin ads
\end{lstlisting}

\begin{teile}
  origin & das Ziel: der Name, unter dem GitHub als Remote eingetragen ist \\
  ads    & der Branch, der hochgeladen wird \\
  -u     & kurz für \cmd{--set-upstream}: verknüpft den lokalen Branch dauerhaft mit dem auf GitHub \\
\end{teile}

Nach dem ersten Push mit \cmd{-u} genügt \cmd{git push} ohne weitere Angaben. Hochgeladen
werden ausschliesslich \textbf{Commits}; Vorgemerktes und Ungespeichertes bleibt lokal.

\subsection*{Die Ausgabe lesen}

\begin{center}
\renewcommand{\arraystretch}{1.25}
\begin{tabularx}{\linewidth}{@{}>{\raggedright\arraybackslash\ttfamily\small}p{5.6cm}>{\raggedright\arraybackslash}X@{}}
  \toprule
  \normalfont\normalsize\textbf{Zeile} & \textbf{Bedeutung} \\
  \midrule
  Enumerating objects: 35 & Git ermittelt, welche gespeicherten Objekte zum Branch gehören \\
  Delta compression using up to 12 threads & überträgt nur Unterschiede statt ganzer Dateien, verteilt auf die Prozessorkerne \\
  Writing objects: 100\% (18/18), 2.29 KiB & tatsächlich übertragen: 18 Objekte, 2,29\,KB \\
  remote: \ldots & Zeilen mit \cmd{remote:} stammen von GitHub, nicht vom eigenen Computer \\
  {*} [new branch] ads -> ads & der Branch existiert jetzt auch auf GitHub \\
  branch \textquotesingle{}ads\textquotesingle{} set up to track \textquotesingle{}origin/ads\textquotesingle{} & Wirkung von \cmd{-u}: lokal und GitHub sind verknüpft \\
  \bottomrule
\end{tabularx}
\end{center}

% =====================================================================
\section{\texorpdfstring{\cmd{git fetch} und \cmd{git pull}}{git fetch und git pull}: Änderungen holen}

\textbf{Zweck:} holt Commits vom Remote.
\textbf{Weg:} GitHub $\rightarrow$ lokales Repository (und Arbeitsordner).

\begin{lstlisting}[style=befehle]
git fetch                              # nur herunterladen, nichts verändern
git pull                               # herunterladen und einbauen
\end{lstlisting}

\cmd{git pull} besteht aus zwei Schritten: Zuerst \cmd{git fetch}, das die neuen Commits
herunterlädt und \cmd{origin/main} aktualisiert, ohne die eigenen Dateien anzurühren. Danach
\cmd{git merge}, das die neuen Commits in den aktuellen Branch einbaut.

\begin{erfahrung}
\cmd{git fetch} ist immer ungefährlich. Wer unsicher ist, holt zuerst mit \cmd{fetch} und
schaut mit \cmd{git status}, ob der eigene Branch vor (\emph{ahead}) oder hinter
(\emph{behind}) GitHub liegt, bevor er mit \cmd{pull} etwas verändert.
\end{erfahrung}

% =====================================================================
\befehl{git stash}{Arbeit kurz beiseitelegen}

\textbf{Zweck:} räumt nicht committete Änderungen vorübergehend weg, damit der Arbeitsordner
sauber ist, zum Beispiel um kurz den Branch zu wechseln.

\begin{lstlisting}[style=befehle]
git stash                              # Änderungen beiseitelegen
git switch main                        # etwas anderes erledigen
git switch ads
git stash pop                          # Änderungen zurückholen
\end{lstlisting}

\begin{zusammenfassung}
\begin{itemize}[nosep]
  \item \cmd{status} und \cmd{diff} zeigen, verändern aber nichts.
  \item \cmd{add}, \cmd{commit}, \cmd{push} bewegen Änderungen nach rechts im Modell,
        \cmd{restore} und \cmd{pull} nach links.
  \item Optionen sind unabhängige Bausteine: einzeln ausprobieren, dann kombinieren.
\end{itemize}
\end{zusammenfassung}
